\section{九大 Skills 类别}

Anthropic 在梳理了内部所有 Skills 之后发现，它们大致聚集为几个反复出现的类别。\textbf{最好的 Skills 清晰地落在某一个类别里；让人困惑的 Skills 往往横跨了好几个类别。}这并非一份穷尽性清单，但它是一个很好的思考框架，帮助你判断组织内是否有遗漏的 Skill 类型。

\subsection{类别一：Library \& API Reference（库与 API 参考）}

帮助正确使用某个库、命令行工具或 SDK 的 Skills。既可以针对内部库，也可以针对 Claude Code 偶尔犯错的常用库。这类 Skills 通常包含：
\begin{itemize}
    \item 一个文件夹的参考代码片段（reference code snippets）
    \item Claude 在写脚本时需要避免的踩坑点（gotchas）列表
\end{itemize}

\textbf{示例}：
\begin{itemize}
    \item \texttt{billing-lib} --- 内部计费库：边界情况、footguns\footnote{Footgun 是指设计上容易让使用者搬起石头砸自己脚的 API 或功能。}等
    \item \texttt{internal-platform-cli} --- 内部 CLI 封装的每个子命令及使用示例
    \item \texttt{frontend-design} --- 让 Claude 更好地适配你的设计系统
\end{itemize}

\subsection{类别二：Product Verification（产品验证）}

描述如何测试或验证代码是否正常工作的 Skills，通常搭配 Playwright、tmux 等外部工具完成验证。

\begin{knowledgebox}{验证类 Skills 值得重点投入}
验证类 Skills 对于确保 Claude 输出的正确性\textbf{极其有用}。作者明确建议：\textbf{值得安排一个工程师花一整周时间专门打磨你的验证 Skills}。

可以考虑的技术手段包括：
\begin{itemize}
    \item 让 Claude 录制输出过程的视频（record a video of its output），这样你可以回看它到底测了什么
    \item 在每一步强制执行程序化的状态断言（programmatic assertions on state at each step）
    \item 在 Skill 中附带各类辅助脚本来实现上述功能
\end{itemize}
\end{knowledgebox}

\textbf{示例}：
\begin{itemize}
    \item \texttt{signup-flow-driver} --- 在无头浏览器中跑完注册 $\to$ 邮件验证 $\to$ 引导流程（onboarding），每一步都可以插入状态断言的钩子
    \item \texttt{checkout-verifier} --- 用 Stripe 测试卡驱动结账 UI，验证发票确实落在正确的状态
    \item \texttt{tmux-cli-driver} --- 针对需要 TTY 的交互式 CLI 测试
\end{itemize}

\subsection{类别三：Data Fetching \& Analysis（数据获取与分析）}

连接到你的数据和监控体系的 Skills。这类 Skills 可能包含：
\begin{itemize}
    \item 带有凭证的数据获取库
    \item 特定的仪表盘 ID
    \item 常见工作流或数据获取方式的说明
\end{itemize}

\textbf{示例}：
\begin{itemize}
    \item \texttt{funnel-query} --- ``哪些事件 join 在一起能看到注册 $\to$ 激活 $\to$ 付费''，加上存放规范 \texttt{user\_id} 的表
    \item \texttt{cohort-compare} --- 对比两个用户群的留存或转化率，标记统计显著的差异（delta），链接到分群定义
    \item \texttt{grafana} --- 数据源 UID、集群名称、问题到仪表盘的对照表
\end{itemize}

\subsection{类别四：Business Process \& Team Automation（业务流程与团队自动化）}

把重复性工作流自动化为一条命令的 Skills。通常指令本身比较简单，但可能有较复杂的对其他 Skills 或 MCP 的依赖。

\begin{knowledgebox}{日志驱动的一致性}
对于这类 Skills，\textbf{把之前的执行结果保存在日志文件中}可以帮助模型保持一致性，并反思之前的执行情况。例如，每次执行后把结果追加到日志，下次运行时模型就能看到历史上下文。
\end{knowledgebox}

\textbf{示例}：
\begin{itemize}
    \item \texttt{standup-post} --- 汇总任务追踪器、GitHub 活动和之前的 Slack 消息，生成格式化的站会汇报（只包含增量变化，delta-only）
    \item \texttt{create-<ticket-system>-ticket} --- 强制执行 schema（有效枚举值、必填字段），加上创建后的工作流（ping reviewer，在 Slack 中发链接）
    \item \texttt{weekly-recap} --- 已合并 PR + 已关闭工单 + 部署记录 $\to$ 格式化周报
\end{itemize}

\subsection{类别五：Code Scaffolding \& Templates（代码脚手架与模板）}

为代码库中的特定功能生成框架样板代码的 Skills。你可以把这些 Skills 和可组合的脚本结合使用。\textbf{当脚手架有自然语言需求、无法纯靠代码覆盖时}，这类 Skills 特别有用。

\textbf{示例}：
\begin{itemize}
    \item \texttt{new-<framework>-workflow} --- 搭建新的 service/workflow/handler，带上你的 annotations
    \item \texttt{new-migration} --- 数据库迁移文件模板加常见踩坑点
    \item \texttt{create-app} --- 新建内部应用，预配认证（auth）、日志（logging）和部署配置（deploy config）
\end{itemize}

\subsection{类别六：Code Quality \& Review（代码质量与审查）}

在团队内部执行代码质量标准并辅助代码审查的 Skills。可以包含确定性的脚本或工具以保证最大程度的可靠性（maximum robustness）。你可能想把这些 Skills 作为 hooks 的一部分自动运行，或放在 GitHub Action 中执行。

\textbf{示例}：
\begin{itemize}
    \item \texttt{adversarial-review} --- 生成一个全新视角的子智能体（subagent）来挑刺，实施修复，反复迭代直到发现的问题退化为吹毛求疵（nitpicks）
    \item \texttt{code-style} --- 强制执行代码风格，特别是 Claude 默认做不好的那些
    \item \texttt{testing-practices} --- 关于如何写测试以及测试什么的指导
\end{itemize}

\begin{knowledgebox}{Subagent 模式}
\texttt{adversarial-review} 中的 subagent 是指 Claude Code 在执行任务时启动的另一个独立 Claude 实例。这个``全新视角''（fresh-eyes）的实例没有见过这段代码的上下文，从而避免了原实例的思维惯性，能发现被忽视的问题。
\end{knowledgebox}

\subsection{类别七：CI/CD \& Deployment（CI/CD 与部署）}

帮助拉取、推送和部署代码的 Skills。这类 Skills 可能引用其他 Skills 来收集数据。

\textbf{示例}：
\begin{itemize}
    \item \texttt{babysit-pr} --- 监控 PR $\to$ 重试不稳定 CI $\to$ 解决合并冲突 $\to$ 启用 auto-merge
    \item \texttt{deploy-<service>} --- 构建 $\to$ 冒烟测试 $\to$ 渐进式流量切换（gradual traffic rollout）并对比错误率 $\to$ 指标恶化时自动回滚
    \item \texttt{cherry-pick-prod} --- 隔离 worktree $\to$ cherry-pick $\to$ 冲突解决 $\to$ 用模板创建 PR
\end{itemize}

\subsection{类别八：Runbooks（运维手册）}

接收一个现象（Slack 消息、告警、错误签名），引导走完多工具排查流程，最后输出结构化报告。

\textbf{示例}：
\begin{itemize}
    \item \texttt{<service>-debugging} --- 把症状映射到工具和查询模式，覆盖流量最高的服务
    \item \texttt{oncall-runner} --- 拉取告警 $\to$ 检查常见嫌疑（usual suspects） $\to$ 格式化排查结论
    \item \texttt{log-correlator} --- 给定一个 request ID，从所有可能碰过这个请求的系统拉取匹配日志
\end{itemize}

\subsection{类别九：Infrastructure Operations（基础设施运维）}

执行日常维护和运维操作的 Skills --- 其中一些涉及破坏性操作，需要安全护栏（guardrails）。这类 Skills 让工程师在执行关键操作时更容易遵循最佳实践。

\begin{warningbox}{破坏性操作需要安全护栏}
基础设施运维类 Skills 可能涉及孤立资源清理、依赖变更等破坏性操作。务必在 Skill 中设计确认流程（如浸泡期 + 用户确认）和安全检查。
\end{warningbox}

\textbf{示例}：
\begin{itemize}
    \item \texttt{<resource>-orphans} --- 找到孤立 Pod/Volume $\to$ 发到 Slack $\to$ 浸泡期（soak period） $\to$ 用户确认 $\to$ 级联清理
    \item \texttt{dependency-management} --- 组织的依赖审批工作流
    \item \texttt{cost-investigation} --- ``为什么我们的存储/出口带宽费用突然飙升''，附带具体的 bucket 和查询模式
\end{itemize}

\subsection{本章小结}

九大类别覆盖了从开发到运维的全生命周期：Library \& API Reference、Product Verification、Data Fetching \& Analysis、Business Process Automation、Code Scaffolding、Code Quality \& Review、CI/CD \& Deployment、Runbooks、Infrastructure Operations。好的 Skill 应当清晰地归属于某一个类别；横跨多个类别的 Skill 往往令人困惑。这个分类框架也可以帮助你审视组织内是否有尚未被覆盖的 Skill 类型。

